\input{sol_pre.tex}
\input{q_2025-12-28.tex}
\begin{document}
\input{sol_head.tex}

\Savol{1}
$\underbrace{22\ldots2}_{50}:22=\underbrace{11\ldots1}_{50}:11$. \ $1111:11=101$, $111111:11=10101$, … juft sondagi birlar $11$ ga bölinganda «$10$» takrorlanib, oxiri $1$ bilan tugaydi: $50$ ta birdan $49$ xonali $\underbrace{1010\ldots10}_{24\text{ ta }10}1$.
\Javob{D}

\Savol{2}
$a=2{,}141414\ldots$, \ $b=2{,}144444\ldots$, \ $c=2{,}114444\ldots$ \ $\Rightarrow$ \ $b>a>c$.
\Javob{B}

\Savol{3}
Jami $1680$; qolgan mahsulot $4$ ga bölinişi kerak, demak sotilgani ham $4$ ga bölinadi: $300$. Qolgan $1380$ dan tarvuz $345$, qovun $315+330+390=1035$.
\Javob{C}

\Savol{4}
Köylak $x$, şim $y$: \ $15x+20y=950\,000$, \ $0{,}72\cdot15x+0{,}8\cdot20y=724\,000$.\\
Birinçisini $0{,}8$ ga köpaytirib ayiramiz: $(12-10{,}8)\cdot x=760\,000-724\,000$ $\Rightarrow$ $1{,}2x=36\,000$, $x=30\,000$.
\Javob{D}

\Savol{5}
$a=100$, $b=14$, $a+b=114$: \ $a^4+b^4+(a+b)^4=2\left(a^2+ab+b^2\right)^2$. Demak
\[\frac{a^3-b^3}{a^2+ab+b^2}=a-b=86.\]
\Javob{A}

\Savol{6}
$100+121+144+169+196=730$, \ $730:146=5$.
\Javob{C}

\Savol{7}
$t=x-2$: surat $t^4+t^2+2t=t(t+1)(t^2-t+2)$, maxraj $t^2-t+2$ $\Rightarrow$ $t(t+1)=(x-2)(x-1)=x^2-3x+2$.
\Javob{A}

\Savol{8}
$a_1+2d=2a_1+d$ $\Rightarrow$ $a_1=d$; hadlar $d,2d,3d$: $6d^3=384$, $d=4$. \ $S_{10}=\frac{2\cdot4+9\cdot4}2\cdot10=220$.
\Javob{A}

\Savol{9}
$b_1b_{19}=b_2b_{18}=\ldots=b_{10}^2$ $\Rightarrow$ $b_1b_2\cdots b_{19}=b_{10}^{19}=5^{19}$.
\Javob{B}

\Savol{10}
$111111^2=12345654321$, \ $666666=6\cdot111111$ $\Rightarrow$ $\sqrt{36}=6$.
\Javob{A}

\Savol{11}
$x\ne0$, $\frac{|x|}{|x|}=1$: \ $\frac{x^3}{|x|}+10x-11=0$.\\
$x>0$: $x^2+10x-11=0$ $\Rightarrow$ $x=1$. \ $x<0$: $x^2-10x+11=0$ $\Rightarrow$ $x=5\pm\sqrt{14}>0$ — mos emas. $1$ ta.
\Javob{A}

\Savol{12}
$14^{\sin x}=7^{\sin x}\cdot2^{\sin x}$ $\Rightarrow$ $2^{\sin x}=2^{-\cos x}$ $\Rightarrow$ $\tg x=-1$, $x=-\frac\pi4+\pi k$.\\
$\left[-4\pi;-\frac{5\pi}2\right]$ da faqat $x=-\frac{13\pi}4$.
\Javob{A}

\Savol{13}
$\arcsin\left(-\frac{\sqrt3}2\right)=-\frac\pi3$ $\Rightarrow$ $\sin2x=-\frac\pi3<-1$ — yeçim yöq.
\Javob{D}

\Savol{14}
$t=\log_3x$: \ $\log_x3=\frac1t$, \ $\log_32\cdot\log_{32}x=\frac{\log_3x}{5}=\frac t5$, \ $\log_{\sqrt x}3=\frac2t$:
\[\frac23-\frac1t=\frac t5-\frac4{5t}\ \Rightarrow\ 3t^2-10t+3=0\ \Rightarrow\ t=3,\ t=\frac13.\]
$x=27$, $x=\sqrt[3]3$; köpaytma $27\sqrt[3]3$.
\Javob{B}

\Savol{15}
$P(x)=3(3x+4)^{20}=\sum C_kx^{20-k}$. \ $P''(x)=\sum C_k(20-k)(19-k)x^{18-k}$; \ $x=-1$ da aynan berilgan ifoda hosil böladi.
\[P''(x)=3\cdot20\cdot19\cdot9\,(3x+4)^{18}\ \Rightarrow\ P''(-1)=10260.\]
\Javob{B}

\Savol{16}
$\alpha+\beta=45^\circ$ bölsa $(1+\tg\alpha)(1+\tg\beta)=2$. \ $7^\circ+38^\circ=8^\circ+37^\circ=45^\circ$ $\Rightarrow$ $2\cdot2=4$.
\Javob{A}

\Savol{17}
$\dfrac{(|x|-2)(|x|-5)}{(x-3)^2}<0$ $\Rightarrow$ $2<|x|<5$, $x\ne3$: \ $-4,-3,4$ — $3$ ta.
\Javob{C}

\Savol{18}
$x_0=3a>0$; \ $y_0=10a^2-a-6-9a^2=a^2-a-6<0$ $\Rightarrow$ $-2<a<3$. \ $a\in\{1;2\}$, yiğindi $3$.
\Javob{C}

\Savol{19}
$xy-x+2y-2=3$ $\Rightarrow$ $(x+2)(y-1)=3$. Natural: $x+2=3$, $y-1=1$ $\Rightarrow$ $x=1$, $y=2$, $x+y=3$.
\Javob{B}

\Savol{20}
Kasr musbat, faqat aniqlaniş sohasi muhim: $\frac{2-x}x\ge0$ $\Rightarrow$ $0<x\le2$, va maxraj $\ne0$.\\
$x=1$: $\sqrt1-\frac22=0$ — mos emas; \ $x=2$: $0-\frac34\ne0$ — mos. $1$ ta.
\Javob{B}

\Savol{21}
$(y-1)x^2+3(y+1)x+4(y+1)=0$ yeçimga ega: $9(y+1)^2-16(y^2-1)\ge0$ $\Rightarrow$ $(y+1)(25-7y)\ge0$ $\Rightarrow$ $y\in\left[-1;\frac{25}7\right]$ ($y=1$ da ham $x=-\frac43$).
\Javob{B}

\Savol{22}
$x^3+(a+3)x^2+(3a+b)x+3b$: \ $a+3=4$, $3b=6$ $\Rightarrow$ $a=1$, $b=2$ ($3a+b=5$ — töğri).
\Javob{A}

\Savol{23}
$\left(\dfrac1{2\cos^2x}\right)'=\dfrac{\sin x}{\cos^3x}$: \quad $\dfrac1{2\cos^2x}\Big|_0^{\pi/4}=1-\dfrac12=\dfrac12$.
\Javob{A}

\Savol{24}
\fig{\begin{tikzpicture}[scale=0.42,font=\small]
\coordinate (A) at (0,0); \coordinate (C) at (15,0); \coordinate (B) at (3.5,6.062);
\draw (A)--(B)--(C)--cycle; \draw (4.5,2.598) circle (2.598); \draw (1.5,0.866) circle (0.866);
\draw[ax] (A)--(4.5,2.598); \draw[axs] (4.5,2.598)--(4.5,0); \draw[axs] (1.5,0.866)--(1.5,0);
\fill (4.5,2.598) circle (4pt); \fill (1.5,0.866) circle (4pt);
\node[below left] at (A) {$A$}; \node[above] at (B) {$B$}; \node[below right] at (C) {$C$};
\node[above right] at (4.5,2.598) {$O_1$}; \node[above left] at (1.5,0.9) {$O_2$};
\end{tikzpicture}}
$\cos A=\frac{49+225-169}{2\cdot7\cdot15}=\frac12$ $\Rightarrow$ $A=60^\circ$. \ $S=\frac12\cdot7\cdot15\cdot\frac{\sqrt3}2=\frac{105\sqrt3}4$, $p=17{,}5$, $r=\frac Sp=\frac{3\sqrt3}2$.\\
Ikkala markaz $A$ bissektrisasida ($\angle=30^\circ$): $AO_1=\frac r{\sin30^\circ}=2r$, $AO_2=2\rho$. Aylanalar urinadi: $AO_1-AO_2=r+\rho$ $\Rightarrow$ $2r-2\rho=r+\rho$ $\Rightarrow$ $\rho=\frac r3=\frac{\sqrt3}2$.
\Javob{C}

\Savol{25}
Bir burçak qolgan ikkitasining yiğindisiga teng $\Rightarrow$ u $90^\circ$; teng yonli töğri burçakli uçburçak. Gipotenuzaga mediana $6$ $\Rightarrow$ gipotenuza $12$, katetlar $6\sqrt2$. \ $S=\frac12\cdot72=36$.
\Javob{A}

\Savol{26}
$\angle ABE=45^\circ$ — içki çizilgan, $\angle AO_1E=90^\circ$ ($AE$ yoyi $90^\circ$).\\
$\triangle AO_1C$ teng yonli: $\angle AO_1C=180^\circ-2\cdot15^\circ=150^\circ$ ($AC$ yoyi $150^\circ$). \ $CE$ yoyi $150^\circ-90^\circ=60^\circ$ $\Rightarrow$ $CE=R_1=6$.\\
Kiçik aylanada $\angle CDE=60^\circ$ $\Rightarrow$ $\angle CO_2E=120^\circ$, $CE=R_2\sqrt3$ $\Rightarrow$ $R_2=2\sqrt3$.\\
$O_1O_2\perp CE$ va uni teng ikkiga böladi: $O_1$ dan $CE$ gaça $6\cos30^\circ=3\sqrt3$, $O_2$ dan $R_2\cos60^\circ=\sqrt3$; markazlar $CE$ ning ikki tomonida:
\[O_1O_2=3\sqrt3+\sqrt3=4\sqrt3.\]
\Javob{C}

\Savol{27}
Qarama-qarşi tomonlar parallel, ular orasidagi masofa $2h$ ($h$ — apofema). $P$ dan har bir juft parallel tomongaça masofalar yiğindisi $2h$: $3\cdot2h=18$ $\Rightarrow$ $h=3$.\\
$h=\frac{\sqrt3}2a$ $\Rightarrow$ $a=2\sqrt3$, \ $S=\frac{3\sqrt3}2a^2=18\sqrt3$.
\Javob{B}

\Savol{28}
$\overrightarrow{AE}+\overrightarrow{EC}=\overrightarrow{AC}$; \ $\overrightarrow{AC}+\overrightarrow{BD}=2\overrightarrow{AD}$, \ $\overrightarrow{AC}-\overrightarrow{BD}=2\overrightarrow{AB}$ $\Rightarrow$ $AD=8$, $AB=5$.
\[S_{AED}=\frac12S_{ABCD}=\frac12\cdot8\cdot5\cdot\sin60^\circ=10\sqrt3.\]
\Javob{B}

\Savol{29}
$AD=a$, $AB=b$: \ $\dfrac{a^2+b^2+ab}{a^2+b^2-ab}=\dfrac{19}7$ $\Rightarrow$ $6a^2-13ab+6b^2=0$ $\Rightarrow$ $\dfrac ab=\dfrac32$.
\Javob{A}

\Savol{30}
\fig{\begin{tikzpicture}[scale=0.5,font=\small]
\coordinate (H) at (0,0); \coordinate (A) at (-4.8,-3.4); \coordinate (B) at (1.5,-4.5); \coordinate (C) at (5.6,-1.2); \coordinate (T) at (0,3.2);
\fill[okgreen!15] (H)--(A)--(B)--(C)--cycle;
\draw (H)--(A)--(B)--(C)--cycle; \draw (T)--(A); \draw (T)--(B); \draw (T)--(C); \draw[axs] (T)--(H); \draw[ax] (H)--(B);
\node[above] at (T) {$T$}; \node[above left] at (H) {$H$}; \node[left] at (A) {$A$}; \node[below] at (B) {$B$}; \node[right] at (C) {$C$};
\end{tikzpicture}}
Oğmalar tekislik bilan $30^\circ$ hosil qiladi: $HA=HB=HC=TH\ctg30^\circ=4\sqrt3$.\\
$\angle AHB+\angle BHC=150^\circ=\angle AHC$, ya'ni $HB$ nur $\angle AHC$ içida. Piramida asosi — $HABC$ törtburçak:
\[S_{HABC}=S_{AHB}+S_{BHC}=\frac12\cdot48\cdot\sin90^\circ+\frac12\cdot48\cdot\sin60^\circ=24+12\sqrt3.\]
$V=\frac13\left(24+12\sqrt3\right)\cdot4=32+16\sqrt3$.
\Javob{B}

\Savol{31}
$x-2$ son $4$ ga ham, $6$ ga ham bölinadi $\Leftrightarrow$ $12$ ga bölinadi: $x=12k+2$.
\Javob{A}

\Savol{32}
Jami: birinçi raqam $4$ xil (nol emas), keyingilari $4$ va $3$ xil: $4\cdot4\cdot3=48$.\\
$4$ ga böliniş: oxirgi ikki raqam $\overline{ab}$ uçun $2a+b\equiv0\pmod4$, $a$ juft $\Rightarrow$ $b\in\{0;4;8\}$.\\
$b=0$: $a$ — $4$ xil, birinçi raqam — $3$ xil: $12$. \ $b=4$: $a=0$ da $3$ ta, $a\in\{2;6;8\}$ da $3\cdot2=6$ ta: $9$. \ $b=8$: $9$.\\
$P=\frac{12+9+9}{48}=\frac{30}{48}=\frac58$.
\Javob{C}

\newpage
\Savol{33}
$\frac43\pi r^3=8\sqrt6\pi$ $\Rightarrow$ $r^3=6\sqrt6$, $r=\sqrt6$. Muntazam oktaedrda (qirrasi $a$) içki şar radiusi $r=\frac{a\sqrt6}6$ $\Rightarrow$ $a=6$.\\
(Oktaedr — asoslari umumiy ikkita muntazam törtburçakli piramida; markazdan yoqqaça masofa: $\frac{a/2\cdot a/\sqrt2}{\sqrt{a^2/4+a^2/2}}=\frac a{\sqrt6}$.)

\textbf{33.} $4\pi r^2=24\pi\approx72$ — \textbf{A}.

\textbf{34.} $8\cdot\frac{\sqrt3}4a^2=2\sqrt3\cdot36=72\sqrt3$ — \textbf{C}.

\textbf{35.} $V=\frac{\sqrt2}3a^3=\frac{\sqrt2}3\cdot216=72\sqrt2$ — \textbf{B}.
\Javob{33 — A, \ 34 — C, \ 35 — B}

\newpage
\Savol{36}
$x^2+xy+y^2=(x+y)^2-xy=\left(x+y-\sqrt{xy}\right)\left(x+y+\sqrt{xy}\right)$ ($x,y\ge0$).\\
Demak $x+y-\sqrt{xy}=\frac{133}{19}=7$. Ikkinçi tenglama bilan: $x+y=13$, $\sqrt{xy}=6$, $xy=36$.\\
$x,y$ — $t^2-13t+36=0$ ildizlari: $4$ va $9$. Yeçimlar: $(4;9)$, $(9;4)$.

a) $\frac{y}{x}$: $\frac94$ va $\frac49$; eng kiçigi $\frac49$.

b) $2$ ta juftlik.
\Javob{a) $\frac49$; \quad b) 2}

\newpage
\Savol{37}
$\sin^3x=\frac{3\sin x-\sin3x}4$, \ $\cos^3x=\frac{3\cos x+\cos3x}4$:
\[\sin^3x\sin3x+\cos^3x\cos3x=\frac{3(\cos x\cos3x+\sin x\sin3x)+\cos^23x-\sin^23x}4=\frac{3\cos2x+\cos6x}4=\cos^32x.\]
Tenglama: $\cos^32x=\cos^36x$ $\Rightarrow$ $\cos2x=\cos6x$ $\Rightarrow$ $6x=\pm2x+2\pi k$ $\Rightarrow$ $x=\frac{\pi k}2$ yoki $x=\frac{\pi k}4$, ya'ni $x=\frac{\pi k}4$.

a) $[-\pi;\pi]$: $k=-4,\dots,4$ — $9$ ta.

b) $(0;2\pi)$: $k=1,\dots,7$, yiğindi $\frac\pi4(1+2+\dots+7)=7\pi$.
\Javob{a) 9; \quad b) $7\pi$}

\newpage
\Savol{38}
$g$ grafigi $Ox$ ga $x=0$ da urinadi: $g(x)=x^2(x-m)$. $f$ grafigi $x=2024$ da urinadi: $f(x)=(x-2024)^2(x-n)$.\\
Kesişiş nuqtalari $Ox$ da yotadi ($g(0)=0$, $f(2024)=0$): $f(0)=0$ $\Rightarrow$ $n=0$; \ $g(2024)=0$ $\Rightarrow$ $m=2024$.
\[g(x)=x^3-2024x^2,\qquad f(x)=x(x-2024)^2=x^3-4048x^2+2024^2x.\]
$p=-2024$, $q=d=0$; \ $a=-4048$, $b=2024^2$, $c=0$.

a) $\dfrac{ap+c}{b+q+d}=\dfrac{4048\cdot2024}{2024^2}=2$.

b) $\dfrac{g(x)}{f(x)}=\dfrac{x^2(x-2024)}{x(x-2024)^2}=\dfrac x{x-2024}$: \ $\dfrac{2025}1+\dfrac{2023}{-1}=2$.
\Javob{a) 2; \quad b) 2}

\newpage
\Savol{39}
Juft-toqlikni tekşiramiz:\\
$h(x)=\dfrac x{1-2^x}-\dfrac x2=\dfrac x2\cdot\dfrac{1+2^x}{1-2^x}$; \ $\dfrac{1+2^{-x}}{1-2^{-x}}=\dfrac{2^x+1}{2^x-1}=-\dfrac{1+2^x}{1-2^x}$ — toq, $x$ — toq $\Rightarrow$ $h$ juft.\\
$\cos^{2025}x$ — juft, $\sin^{2024}x$ — juft, $\ln\sqrt{x+\sqrt{1+x^2}}$ — toq ($(-x+\sqrt{1+x^2})(x+\sqrt{1+x^2})=1$).\\
Demak $f$ — toq. Toq funksiya hosilasi juft, juftniki — toq: $f'$ juft, $f''$ toq, $f'''$ juft.

a) $f'(2025)-f'(-2025)=0$, maxraj $2f'(1)$ $\Rightarrow$ $0$.

b) $f''(2)+f''(-2)=0$, maxraj $2f'''(1)$ $\Rightarrow$ $0$.
\Javob{a) 0; \quad b) 0}

\newpage
\Savol{40}
Aniqlaniş sohasi $[2;6]$: $ax^2+bx+c=a(x-2)(x-6)$, $a<0$. \ $f(2)=d=4$. Eng katta qiymat $x=4$ da: $\sqrt{-4a}+4=6$ $\Rightarrow$ $a=-1$:
\[f(x)=\sqrt{-x^2+8x-12}+4,\quad a=-1,\ b=8,\ c=-12,\ d=4.\]
$y=f(x)$ — markazi $(4;4)$, radiusi $2$ yarim aylana. \ $g$ $Ox$ ga $C(4;0)$ da urinadi: $g=k(x-4)^2$, $g(2)=4$ $\Rightarrow$ $k=1$: $g=x^2-8x+16$.

a) $\dfrac{k+p+q}{a+b+c+d}=\dfrac{1-8+16}{-1+8-12+4}=\dfrac9{-1}=-9$.

b) \[S=\int_2^6\bigl(f(x)-g(x)\bigr)dx=\underbrace{\frac{\pi\cdot2^2}2}_{\text{yarim doira}}+4\cdot4-\int_2^6(x-4)^2dx=2\pi+16-\frac{16}3=2\pi+\frac{32}3.\]
\Javob{a) $-9$; \quad b) $2\pi+\frac{32}3$}

\newpage
\Savol{41}
\fig{\begin{tikzpicture}[scale=0.38,font=\small]
\coordinate (A) at (0,0); \coordinate (C) at (14,0); \coordinate (B) at (5,12); \coordinate (O) at (6,4); \coordinate (E) at (1.667,4); \coordinate (F) at (11,4);
\fill[okgreen!12] (B)--(E)--(F)--cycle;
\draw (A)--(B)--(C)--cycle; \draw (O) circle (4); \draw (E)--(F); \draw[axs] (O)--(6,0); \draw[ax] (B)--(5,0); \fill (O) circle (4pt);
\node[below left] at (A) {$A$}; \node[above] at (B) {$B$}; \node[below right] at (C) {$C$}; \node[left] at (E) {$E$}; \node[right] at (F) {$F$}; \node[above right] at (O) {$O$};
\node[aux,right] at (6,2) {$r$}; \node[aux,left] at (5,6.5) {$h$};
\end{tikzpicture}}
$S_{ABC}=37\frac13+46\frac23=84$. \ $\triangle BEF\sim\triangle BAC$: $k^2=\frac{112/3}{84}=\frac49$, $k=\frac23$.\\
Balandlik $h=\frac{2\cdot84}{14}=12$. $EF$ dan $AC$ gaça masofa $r$ (markazdan asosgaça), $B$ dan $EF$ gaça $kh$.

a) $r=h-kh=\frac13\cdot12=4$.

b) $p=\frac Sr=21$ $\Rightarrow$ $AB+BC=42-14=28$. \ $AE=\frac13AB$, $CF=\frac13BC$, $EF=\frac23\cdot14=\frac{28}3$:
\[P_{AEFC}=\frac{28}3+\frac{28}3+14=\frac{98}3.\]
\Javob{a) 4; \quad b) $\frac{98}3$}

\newpage
\Savol{42}
\fig{\begin{tikzpicture}[scale=0.75,font=\small]
\coordinate (A) at (0.000,0.000); \coordinate (B) at (3.247,3.802); \coordinate (C) at (5.843,3.802); \coordinate (D) at (14.000,0.000); \coordinate (E) at (5.000,0.000); \coordinate (G) at (4.124,1.901); \coordinate (H) at (5.421,1.901); \coordinate (P) at (4.545,2.095); \coordinate (M) at (1.624,1.901); \coordinate (N) at (9.921,1.901);
\draw (A)--(B)--(C)--(D)--cycle; \draw (B)--(E)--(C); \draw (A)--(P)--(D); \draw[ax] (M)--(N);
\draw[axs] (P)--(E);
\node[below left] at (A) {$A$}; \node[above] at (B) {$B$}; \node[above] at (C) {$C$}; \node[below right] at (D) {$D$}; \node[below] at (E) {$E$};
\node[below left] at (G) {$G$}; \node[below right] at (H) {$H$}; \node[above] at (P) {$P$}; \node[left] at (M) {$M$}; \node[above right] at (N) {$N$};
\end{tikzpicture}}
$\angle AEB=\angle EBC=\angle ABE$ $\Rightarrow$ $AE=AB=5$; xuddi şunday $DE=DC=9$.\\
$\triangle ABE$ teng yonli, $AG$ — bissektrisa $\Rightarrow$ $AG\perp BE$, $G$ — $BE$ örtasi. $\triangle DCE$ da $DH\perp CE$, $H$ — $CE$ örtasi.

a) $G$ va $H$ — $BE$ va $CE$ örtalari, ular trapetsiya örta çiziğida yotadi. Demak $MN$ — örta çiziq, $N$ — $CD$ örtasi: $\frac{CN}{ND}=1$.

b) $u=\frac{\angle A}2$, $v=\frac{\angle D}2$. \ $GE=5\sin u$, $HE=9\sin v$. Trapetsiya balandligi: $h=BE\sin\angle BEA=10\sin u\cos u$, şuningdek $h=18\sin v\cos v$:
\[5\sin u\cos u=9\sin v\cos v.\qquad(*)\]
$\angle PGE=\angle PHE=90^\circ$ $\Rightarrow$ $P,G,E,H$ — $PE$ diametrli aylanada. \ $\angle HPE=\angle HGE=\angle GEA=90^\circ-u$ ($GH\parallel AD$), \ $\angle GPE=\angle GHE=90^\circ-v$:
\[PH=HE\tg u,\qquad PG=GE\tg v.\]
\[\frac{PH}{GE}=\frac{9\sin v\tg u}{5\sin u}=\frac{9\sin v}{5\cos u}\overset{(*)}=\frac{\sin u}{\cos v},\qquad \frac{GP}{HE}=\frac{5\sin u\tg v}{9\sin v}=\frac59\cdot\frac{\sin u}{\cos v}=\frac59\cdot\frac{PH}{GE}=\frac59\cdot\frac37=\frac5{21}.\]
\Javob{a) 1; \quad b) $\frac5{21}$}

\newpage
\Savol{43}
\fig{\begin{tikzpicture}[scale=0.7,font=\small]
\coordinate (O) at (0,0); \coordinate (A) at (-3.449,-0.318); \coordinate (C) at (3.449,-0.318); \coordinate (B) at (-1.449,3.146);
\coordinate (K) at ($(A)!0.5!(B)$); \coordinate (F) at ($(B)!0.5!(C)$); \coordinate (M) at (-3.162,1.414); \coordinate (N) at (3.162,1.414);
\draw (O) circle (3.464); \draw (A)--(B)--(C)--cycle; \draw (M)--(N);
\draw[axs] (O)--(F); \draw[ax] (O)--(0,1.414); \draw[ax] (O)--(0,-0.318); \draw[axs] (O)--(N);
\fill (O) circle (2pt);
\node[left] at (A) {$A$}; \node[above] at (B) {$B$}; \node[right] at (C) {$C$}; \node[left] at (M) {$M$}; \node[right] at (N) {$N$};
\node[below right] at (K) {$K$}; \node[above right] at (F) {$F$}; \node[below left] at (O) {$O$}; \node[above] at (0,1.414) {$L$};
\end{tikzpicture}}
$BC=2R\sin A=4\sqrt3\cdot\frac{\sqrt3}2=6$ — mos.

a) $OF\perp BC$ ($F$ — vatar örtasi): $OF=\sqrt{R^2-\left(\frac{BC}2\right)^2}=\sqrt{12-9}=\sqrt3$.

b) $\triangle ABC$: $36=16+AC^2-4AC$ $\Rightarrow$ $AC=2+2\sqrt6$. \ $KF\parallel AC$ (örta çiziq) va $AC$ dan $\frac12h_B=\frac12\cdot4\sin60^\circ=\sqrt3$ masofada.\\
$O$ dan $AC$ gaça: $\sqrt{R^2-\left(\frac{AC}2\right)^2}=\sqrt{12-(1+\sqrt6)^2}=\sqrt{5-2\sqrt6}=\sqrt3-\sqrt2$ ($\angle B$ ötkir — $O$ va $B$ bir tomonda).\\
$O$ dan $MN$ gaça: $OL=\sqrt3-(\sqrt3-\sqrt2)=\sqrt2$.
\[MN=2\sqrt{R^2-OL^2}=2\sqrt{12-2}=2\sqrt{10}.\]
\Javob{a) $\sqrt3$; \quad b) $2\sqrt{10}$}

\newpage
\Savol{44}
\fig{\begin{tikzpicture}[scale=0.62,font=\small]
\coordinate (P1) at (0,0); \coordinate (P2) at (7,0); \coordinate (P3) at (3.5,6.062);
\coordinate (B) at (1,0); \coordinate (C) at (5,0); \coordinate (D) at ($(P2)!2/7!(P3)$); \coordinate (E) at ($(P2)!6/7!(P3)$); \coordinate (F) at ($(P1)!6/7!(P3)$); \coordinate (A) at ($(P1)!1/7!(P3)$);
\fill[okgreen!12] (A)--(B)--(C)--(D)--(E)--(F)--cycle;
\draw[ax] (P1)--(P2)--(P3)--cycle; \draw (A)--(B)--(C)--(D)--(E)--(F)--cycle;
\node[below] at (B) {$B$}; \node[below] at (C) {$C$}; \node[right] at (D) {$D$}; \node[right] at (E) {$E$}; \node[left] at (F) {$F$}; \node[left] at (A) {$A$};
\node[below] at (0.5,0) {\scriptsize$1$}; \node[below] at (3,0) {\scriptsize$4$}; \node[below] at (6,0) {\scriptsize$2$};
\end{tikzpicture}}
Barça burçaklari $120^\circ$. $BC$, $DE$, $FA$ tomonlarni davom ettirsak, burçaklarda muntazam uçburçaklar kesiladi (tomonlari $AB$, $CD$, $EF$) va katta muntazam uçburçak hosil böladi:
\[AB+BC+CD=CD+DE+EF=EF+FA+AB\ \Rightarrow\ 1+4+2=2+4+EF=EF+FA+1.\]
$EF=1$, $FA=5$.

a) $P=1+4+2+4+1+5=17$.

b) Katta uçburçak tomoni $7$: \ $S=\frac{\sqrt3}4\left(7^2-1^2-2^2-1^2\right)=\frac{43\sqrt3}4$.
\[V=\frac13\cdot\frac{43\sqrt3}4\cdot4\sqrt3=43.\]
\Javob{a) 17; \quad b) 43}

\newpage
\Savol{45}
\fig{\begin{tikzpicture}[scale=0.42,font=\small]
\draw[very thick] (-2,0)--(16,0); \draw (0,0)--(0,-8); \draw[dashed] (0,-8)--(15,0);
\draw[axs] (0,-8)--(10.667,0); \draw[ax] (0,-8)--(6,0);
\fill (0,-8) circle (4pt); \fill (15,0) circle (4pt); \fill (10.667,0) circle (3pt); \fill (6,0) circle (3pt);
\node[above] at (0,0) {$D$}; \node[above] at (15,0) {$H$}; \node[below] at (0,-8) {$A$}; \node[above] at (10.667,0) {$X$}; \node[above] at (6,0) {$Y$};
\node[left] at (0,-4) {$8$};
\end{tikzpicture}}
$DH=\sqrt{17^2-8^2}=15$. Maşina dala böylab yölning $DX=x$ nuqtasiga çiqib, qolganini yöldan yuradi.

a) Yuk bilan: $t(x)=\dfrac{\sqrt{64+x^2}}8+\dfrac{15-x}{10}$, \ $t'(x)=\dfrac x{8\sqrt{64+x^2}}-\dfrac1{10}=0$ $\Rightarrow$ $10x=8\sqrt{64+x^2}$ $\Rightarrow$ $x=\dfrac{32}3$, $AX=\dfrac{40}3$.
\[t_{\min}=\frac{40}{3\cdot8}+\frac{15-\frac{32}3}{10}=\frac53+\frac{13}{30}=\frac{63}{30}\text{ soat}=126\text{ daqiqa}.\]
b) Yuksiz: $t(x)=\dfrac{\sqrt{64+x^2}}{12}+\dfrac{15-x}{20}$, \ $t'=0$ $\Rightarrow$ $20x=12\sqrt{64+x^2}$ $\Rightarrow$ $x=6$, $AY=10$.
\[t_{\min}=\frac{10}{12}+\frac9{20}=\frac{77}{60}\text{ soat}=77\text{ daqiqa}.\]
\Javob{a) 126; \quad b) 77}

\ShowJavoblar
\end{document}
